\section{Biological Example Family III: Rhythmic Control and Limit Cycles}

Some systems are not best understood as settling into fixed points at all.
Instead, they generate recurrent oscillatory activity. This is where limit
cycles and central-pattern style reasoning become useful.

\begin{definition}[Limit cycle]
A \textbf{limit cycle} is a stable periodic trajectory. Nearby trajectories are
drawn toward a repeating orbit rather than toward a static equilibrium.
\end{definition}

\begin{discussion}
Rhythmic biological systems such as locomotor control, breathing-like rhythms,
and repetitive patterned action are pedagogically important because they block
a common mistake: not every useful regime is a still basin. Some regimes are
productive precisely because they are cyclic. Once the cycle is entered, the
system continues through repeated phases with low moment-to-moment decision
cost.
\end{discussion}

\begin{example}[Exercise routine as cycle entry]
An exercise routine often fails before the first motion, not during repetition.
That is mathematically suggestive. Entering the regime may require crossing a
large barrier, but once the session is underway, repeating sub-actions can
become rhythmically easy. The practical lesson is that intervention design
should distinguish \emph{cycle entry} from \emph{cycle maintenance}. The two are
not the same control problem.
\end{example}
